% GENERATED FILE. Edit canonical lessons, then run npm run generate:books.
% canonical-source-hash: fnv1a64:d39a9652ff43ae68
% canonical-lessons: SA-C162-vetanam, SA-C162-rnam, SA-C162-labhah, SA-C162-hanih, SA-C162-vyaparah

\chapter{Wages, Debt, Profit, Loss, Trade}
\label{ch:sa-wages-debt-profit-loss-trade}
\hlchaptermodality{162}

\begin{chapteropening}
\textbf{By the end of this chapter:} \emph{I can say wages, a debt, profit, loss and trade.}

Complete the last lesson of chapter 162: I can say wages, a debt, profit, loss and trade.
\end{chapteropening}

\section[\texorpdfstring{\sk{वेतनम्} (vetanam)}{वेतनम् (vetanam)}]{\sk{वेतनम्} (vetanam) — wages, salary}

\label{lesson:SA-C162-vetanam}



\begin{quote}
Before the new one: say the Sanskrit for an anklet, then the Sanskrit for an earring.
\end{quote}

\subsection*{You'll want to know: \sk{वेतनम्}}
\textbf{\sk{वेतनम्}} — \emph{vetanam} — \textquotedblleft{}wages, salary\textquotedblright{}.

\begin{grammarlens}[title={What you've built}]
One more word for this chapter.
\end{grammarlens}

\subsection*{Your turn}
\begin{itemize}
\raggedright
  \item \emph{Say it:} \emph{vetanam}
  \item \emph{Say it:} \emph{vetanam}, once more
  \item \emph{Say it:} say \emph{vetanam}
  \item \emph{Recall:} say the Sanskrit for an anklet, then the Sanskrit for an earring, then say \emph{vetanam} again
\end{itemize}

\begin{tcolorbox}[breakable,colback=teal!4,colframe=teal!35!black,title={Before you move on}]
What does \sk{वेतनम्} mean? (\textquotedblleft{}Wages, salary\textquotedblright{}.) Say it once more.
\end{tcolorbox}
\section[\texorpdfstring{\sk{ऋणम्} (ṛṇam)}{ऋणम् (ṛṇam)}]{\sk{ऋणम्} (ṛṇam) — a debt}

\label{lesson:SA-C162-rnam}



\begin{quote}
Before the new one: say the Sanskrit for an earring, then the Sanskrit for wages.
\end{quote}

\subsection*{You'll want to know: \sk{ऋणम्}}
\textbf{\sk{ऋणम्}} — \emph{ṛṇam} — \textquotedblleft{}a debt\textquotedblright{}.

\begin{grammarlens}[title={What you've built}]
One more word for this chapter.
\end{grammarlens}

\subsection*{Your turn}
\begin{itemize}
\raggedright
  \item \emph{Say it:} \emph{ṛṇam}
  \item \emph{Say it:} \emph{ṛṇam}, once more
  \item \emph{Say it:} say \emph{ṛṇam}
  \item \emph{Recall:} say the Sanskrit for an earring, then the Sanskrit for wages, then say \emph{ṛṇam} again
\end{itemize}

\begin{tcolorbox}[breakable,colback=teal!4,colframe=teal!35!black,title={Before you move on}]
What does \sk{ऋणम्} mean? (\textquotedblleft{}A debt\textquotedblright{}.) Say it once more.
\end{tcolorbox}
\section[\texorpdfstring{\sk{लाभः} (lābhaḥ)}{लाभः (lābhaḥ)}]{\sk{लाभः} (lābhaḥ) — profit, gain}

\label{lesson:SA-C162-labhah}



\begin{quote}
Before the new one: say the Sanskrit for wages, then the Sanskrit for a debt.
\end{quote}

\subsection*{You'll want to know: \sk{लाभः}}
\textbf{\sk{लाभः}} — \emph{lābhaḥ} — \textquotedblleft{}profit, gain\textquotedblright{}.

\begin{grammarlens}[title={What you've built}]
One more word for this chapter.
\end{grammarlens}

\subsection*{Your turn}
\begin{itemize}
\raggedright
  \item \emph{Say it:} \emph{lābhaḥ}
  \item \emph{Say it:} \emph{lābhaḥ}, once more
  \item \emph{Say it:} say \emph{lābhaḥ}
  \item \emph{Recall:} say the Sanskrit for wages, then the Sanskrit for a debt, then say \emph{lābhaḥ} again
\end{itemize}

\begin{tcolorbox}[breakable,colback=teal!4,colframe=teal!35!black,title={Before you move on}]
What does \sk{लाभः} mean? (\textquotedblleft{}Profit, gain\textquotedblright{}.) Say it once more.
\end{tcolorbox}
\section[\texorpdfstring{\sk{हानिः} (hāniḥ)}{हानिः (hāniḥ)}]{\sk{हानिः} (hāniḥ) — loss}

\label{lesson:SA-C162-hanih}



\begin{quote}
Before the new one: say the Sanskrit for a debt, then the Sanskrit for profit.
\end{quote}

\subsection*{You'll want to know: \sk{हानिः}}
\textbf{\sk{हानिः}} — \emph{hāniḥ} — \textquotedblleft{}loss\textquotedblright{}.

\begin{grammarlens}[title={What you've built}]
One more word for this chapter.
\end{grammarlens}

\subsection*{Your turn}
\begin{itemize}
\raggedright
  \item \emph{Say it:} \emph{hāniḥ}
  \item \emph{Say it:} \emph{hāniḥ}, once more
  \item \emph{Say it:} say \emph{hāniḥ}
  \item \emph{Recall:} say the Sanskrit for a debt, then the Sanskrit for profit, then say \emph{hāniḥ} again
\end{itemize}

\begin{tcolorbox}[breakable,colback=teal!4,colframe=teal!35!black,title={Before you move on}]
What does \sk{हानिः} mean? (\textquotedblleft{}Loss\textquotedblright{}.) Say it once more.
\end{tcolorbox}
\section[\texorpdfstring{\sk{व्यापारः} (vyāpāraḥ)}{व्यापारः (vyāpāraḥ)}]{\sk{व्यापारः} (vyāpāraḥ) — trade, business}

\label{lesson:SA-C162-vyaparah}



\begin{quote}
Before the new one: say the Sanskrit for profit, then the Sanskrit for loss.
\end{quote}

\subsection*{You'll want to know: \sk{व्यापारः}}
\textbf{\sk{व्यापारः}} — \emph{vyāpāraḥ} — \textquotedblleft{}trade, business\textquotedblright{}.

\begin{grammarlens}[title={What you've built}]
One more word for this chapter.
\end{grammarlens}

\subsection*{Your turn}
\begin{itemize}
\raggedright
  \item \emph{Say it:} \emph{vyāpāraḥ}
  \item \emph{Say it:} \emph{vyāpāraḥ}, once more
  \item \emph{Say it:} say \emph{vyāpāraḥ}
  \item \emph{Recall:} say the Sanskrit for profit, then the Sanskrit for loss, then say \emph{vyāpāraḥ} again
\end{itemize}

\begin{tcolorbox}[breakable,colback=teal!4,colframe=teal!35!black,title={Before you move on}]
What does \sk{व्यापारः} mean? (\textquotedblleft{}Trade, business\textquotedblright{}.) Say it once more.
\end{tcolorbox}
